where $\mathbf{x}_{\text{MRI}}$ denotes the input magnetic resonance image, $\mathbf{y}_f$ and $\mathbf{y}_a$ are the corresponding ground-truth FDG and AV45 PET images (when available), and $\delta_F, \delta_A \in \{0,1\}$ indicate the presence ($1$) or absence ($0$) of each PET modality for a given subject. The loss weights were $\lambda_F=1.0$ and $\lambda_A=1.0$. This formulation ensures that supervision is applied solely to available modalities, enabling training in heterogeneous clinical cohorts with incomplete multi-tracer imaging.
